\documentclass[10pt,a4paper]{article}
\usepackage{fontspec}
\usepackage{polyglossia}
\setdefaultlanguage{english}
\setmainfont{DejaVu Sans}
\usepackage{geometry}
\geometry{left=1.5cm,right=1.5cm,top=1.5cm,bottom=1.5cm}
\usepackage{enumitem}
\usepackage{titlesec}
\usepackage{xcolor}
\definecolor{linkbtn}{HTML}{1F5FD6}
\usepackage{hyperref}
\hypersetup{colorlinks=true,linkcolor=black,urlcolor=linkbtn,pdftitle={CV - Leonid Nekrylov}}
\titleformat{\section}{\large\bfseries}{}{0em}{}[\titlerule]
\titlespacing{\section}{0pt}{1ex}{1ex}
\setlength{\parindent}{0pt}
\newcommand{\entry}[4]{%
  \par\noindent\textbf{#1}\hfill\textit{#2}\par
  \noindent\textit{#3}\hfill\textit{#4}\par\vspace{2pt}}
\newcommand{\entrygap}{\par\vspace{0.5cm}}
\setlength{\fboxsep}{2.5pt}
\newcommand{\projlink}[3]{%
  \par\noindent\href{#1}{\colorbox{linkbtn}{\textcolor{white}{\small\textbf{▶~#2}}}}%
  \hspace{6pt}\href{#1}{\small\underline{#3}}\par\vspace{3pt}}

\begin{document}
\pagestyle{empty}
\begin{center}
{\LARGE \textbf{Leonid Nekrylov}}
\vspace{0.1cm}

{\large Robot Learning Engineer}
\vspace{0.2cm}

Moscow, Russia |
\href{mailto:nekrylov.ln@phystech.edu}{nekrylov.ln@phystech.edu}
\\
\href{https://t.me/leonid_n_n}{t.me/leonid\_n\_n}
\end{center}

\section*{Education}
\textbf{Moscow Institute of Physics and Technology (MIPT)} \hfill \textit{2023 -- 2027} \\
\textit{BSc, Phystech School of Applied Mathematics and Computer Science} \hfill Dolgoprudny, Russia \\
Major: Applied Mathematics and Physics

\subsection*{Relevant courses}
\vspace{0.1cm}
\begin{itemize}[leftmargin=*,nosep]
\item \textbf{Optimization and ML:} Year-long course on optimization methods in ML (D. Merkulov), Machine Learning 1 (R. Neychev), LLMs (Giraffe AI).
\item \textbf{Reinforcement Learning:} Self-studied UC Berkeley CS285 (\textit{Deep Reinforcement Learning}, Sergey Levine).
\item \textbf{Mathematics:} Probability Theory, Stochastic Processes, Mathematical Statistics, Linear Algebra, Calculus and Functional Analysis.
\item \textbf{Computer Science:} Algorithms and Data Structures, C++, Databases (PostgreSQL), Python.
\item \textbf{Physics/Engineering:} Introduction to Robotics (O. Bulichev), Theoretical Mechanics.
\end{itemize}

\section*{Experience and projects}
\entry{\parbox[t]{0.65\textwidth}{UNITAS: investigating a 3D World-Action Model for robotic manipulation}}{November 2026 -- present}{Research project}{}
\projlink{https://github.com/leonidnn/unitas-wam-evaluation.git}{See the code}{github.com/leonidnn/unitas-wam-evaluation}
\begin{itemize}[leftmargin=*,nosep]
\item Deployed UNITAS, a 1.7B-parameter model for generating robot actions and predicting object motion in 3D.
\item Prepared a reproducible policy evaluation in LIBERO: compared lighting conditions using matched initial states, compared checkpoints, and measured inference latency.
\item Implemented simulation episode recording and local replay in MuJoCo from arbitrary viewpoints without rerunning the neural network.
\item Evaluated model transfer to a custom scene: prepared input from RGB video with LLM-estimated depth and compared projected predicted trajectories against held-out frames.
\end{itemize}
\entrygap

\entry{Pet or Threat, behaviour of a Unitree Go1 quadruped robot}{October 2026}{Personal project, MuJoCo}{Moscow, Russia}
\projlink{https://leonidnn.github.io/pet-or-threat-site/}{Watch the video}{leonidnn.github.io/pet-or-threat-site}
\begin{itemize}[leftmargin=*,nosep]
\item Implemented autonomous behaviour selection for a simulated quadruped robot without external commands: guarding, play and social interaction, waiting for its owner, and rest.
\item The decision module receives perception data and generates high-level movement and gaze commands; pretrained RL policies handle locomotion. The module is simulator-independent; adapting it to another robot requires replacing the low-level controller.
\end{itemize}
\entrygap

\entry{MIPT CCM, RL agent for RGB-D object grasping}{September 2026 -- October 2026}{Research project, Isaac Sim, UR10e}{Moscow, Russia}
\projlink{https://github.com/leonidnn/isaac-rgbd-grasping}{Watch the video and view the code}{github.com/leonidnn/isaac-rgbd-grasping}
\begin{itemize}[leftmargin=*,nosep]
\item Built an Isaac Sim environment for grasp training with object randomization. Developed a custom Vulkan layer for explicit GPU selection without using other devices on a shared multi-user server.
\item Trained an RL agent: a network estimates grasp success for each point and angle from a depth map, and the agent selects the highest-scoring option.
\item Evaluation in new random poses with a full manipulator approach yielded 18/20 and 19/20 successful grasps on two objects. SAC did not learn the task within 5.7k attempts.
\item Evaluated skill transfer to a challenging deformable object (a bag) and compared fine-tuning strategies that preserve performance on previous objects. The best tested approach used two networks switched by object type: 88\% on the new object versus 42\% without fine-tuning, and 95\% on previous objects.
\item Identified a rendering artefact that distorted training and corrected the setup.
\end{itemize}
\entrygap

\entry{NB-Technologies, Hybrid Solver}{September 2025 -- November 2025}{LLM developer}{Moscow, Russia}
\projlink{https://huggingface.co/spaces/leonid-nn/hybrid-solver}{Try the solver}{huggingface.co/spaces/leonid-nn/hybrid-solver}
\projlink{https://github.com/leonidnn/hybrid-solver}{See the code}{github.com/leonidnn/hybrid-solver}
\begin{itemize}[leftmargin=*,nosep]
\item Developed a mathematical problem solver that verifies answers and abstains from returning a result when verification fails.
\item Applied the principle of moving from a Monte Carlo algorithm to a Las Vegas algorithm (an analogy with the ZPP class): an answer is returned only when three independent methods agree~--- language-model reasoning, exact computation in SymPy, and numerical verification.
\end{itemize}
\entrygap

\entry{Cool Deal, wearable climate-control systems}{December 2024 -- September 2025}{Founder, project lead}{Moscow, Russia}
\begin{itemize}[leftmargin=*,nosep]
\item Led a team of 5 spanning engineering and sales from the initial idea to a working prototype of a wearable climate-control system.
\item Designed and built prototypes: CAD, sensors, microcontroller firmware, and bench testing.
\item Secured a 1M rouble ``Student Startup'' grant for R\&D and testing.
\item Handled sales and carried out 2 pilot deployments. Decided to close the project based on the results of hypothesis testing.
\end{itemize}
\entrygap

\entry{Sber (Recommender Systems Team / Amazme)}{July 2024 -- November 2024}{Data Scientist}{Moscow, Russia}
\begin{itemize}[leftmargin=*,nosep]
\item Built PySpark data marts from raw clickstream.
\item Fine-tuned and validated candidate-retrieval models for a recommender system in PyTorch using offline metrics.
\item Traced the causes of recommender-system anomalies from bug reports and developed unit tests.
\end{itemize}
\vspace{0.2cm}

\section*{Technical skills}
\begin{itemize}[leftmargin=*,nosep]
\item Python (NumPy, SciPy, PyTorch, Matplotlib, PySpark, SymPy), PostgreSQL, Git, Pytest, Linux, \LaTeX, Docker, ROS2, NVIDIA Isaac Sim / Isaac Lab, MuJoCo.
\item English~-- C1 (TOEFL iBT 107).
\end{itemize}
\end{document}
